\subsection{Quotient by a regular sequence}

PROPOSITION 4. --- Let A be a Noetherian ring, J an ideal of A generated by an A-regular sequence x, and Ω a finitely generated A-module.

\begin{enumerate}
\def\labelenumi{\alph{enumi})}
\item
  If the A-module \(\Omega\) is dualizing, the sequence \(\mathbf{x}\) is \(\Omega\) -regular and the A/J-module \(\Omega/\mathrm{J}\Omega\) is dualizing;
\item
  If the A/J-module \(\Omega/\mathrm{J}\Omega\) is dualizing, that J is contained in the radical of \(\Lambda\) and that the sequence x is \(\Omega\) -regular, the A-module \(\Omega\) is dualizing.
\end{enumerate}

Reasoning by induction on the length of the sequence x, we reduce to the case where it is reduced to a single element x. Suppose that the A-module \(\Omega\) let it be dualizing. For every maximal ideal m of A containing x, we have \(\dim(\mathrm{A}_{\mathfrak{m}}/x\mathrm{A}_{\mathfrak{m}})=\dim(\mathrm{A}_{\mathfrak{m}})-1\) (VIII, § 3, no. 1, cor. 2), and consequently \(\mathrm{Hom}_{\mathrm{A}_{\mathfrak{m}}}\left(\mathrm{A}_{\mathfrak{m}} / x\mathrm{A}_{\mathfrak{m}}, \Omega_{\mathfrak{m}}\right) = 0\) (no. 1, prop. 3, a)). This implies \(\mathrm{Hom}_{\mathrm{A}}(\mathrm{A} / x\mathrm{A}, \Omega) = 0\) so that the homothety \(x_{\Omega}\) is injective. One may therefore assume, in order to prove the proposition, that the homothety \(x_{\Omega}\) is injective.

Let \(\overline{\mathsf{A}}\) the ring A/xA; let \(\mathfrak{m}\) a maximal ideal of A containing \(x\) , and let \(\overline{\mathfrak{m}}\) its image in \(\overline{\mathsf{A}}\) . The A-module A/ \(\mathfrak{m}\) is annihilated by \(x\) and is identified with \(\overline{\mathsf{A}}/\overline{\mathfrak{m}}\) ; thus for every integer \(i \geqslant 1\) we have an isomorphism \(\operatorname{Ext}_{\mathsf{A}}^{i}(\mathsf{A}/\mathfrak{m}, \Omega) \longrightarrow \operatorname{Ext}_{\overline{\mathsf{A}}}^{i-1}(\overline{\mathsf{A}}/\overline{\mathfrak{m}}, \Omega/x\Omega)\) ( \(\S\) 3, no. 4, prop. 7). We have

\[
\mathrm{ht} (\overline {{\mathfrak {m}}}) = \dim (\overline {{\mathrm{A}}} _ {\overline {{\mathfrak {m}}}}) = \dim (\mathrm{A} _ {\mathfrak {m}} / x \mathrm{A} _ {\mathfrak {m}}) = \dim (\mathrm{A} _ {\mathfrak {m}}) - 1 = \mathrm{ht} (\mathfrak {m}) - 1
\]

(VIII, § 3, no. 1, cor. 2, a)). Now the maximal ideals of \(\overline{\mathrm{A}}\) are the ideals \(\overline{\mathfrak{m}}\) , where \(\mathfrak{m}\) is a maximal ideal of A containing \(x\) if moreover \(x\) belongs to the radical of A, every maximal ideal of A contains \(x\) The proposition follows from this.

COROLLARY 1. Let A be an integral noetherian ring. Every dualizing A-module is torsion-free and of rank 1.

Let \(\Omega\) a dualizing A-module; it is torsion-free by prop. 4. Let K be the field of fractions of A; the K-vector space \(K \otimes_{\Lambda} \Omega\) is dualizing \(n^{\circ}\) 1, prop. 2), hence of dimension 1.

COROLLARY 2. Let A be a local Macaulay ring, Ω a finitely generated A-module, and x a maximal secant sequence of elements of m \(_{A}\) , generating an ideal J. The following conditions are equivalent:

\begin{enumerate}
\def\labelenumi{(\roman{enumi})}
\item
  the A-module Ω is dualizing;
\item
  the A-module \(\Omega\) is Macaulay of dimension equal to \(\dim(A)\) , and by \(\Omega/J\Omega\) is an indecomposable injective module over the local artinian ring \(A/J\) ;
\item
  the sequence x is Ω-regular and Ω/JΩ is an indecomposable injective module over the local artinian ring A/J ;
\item
  the sequence \(\mathbf{x}\) is \(\Omega\) -regular, one has \(\mathrm{long}_{\mathrm{A}}(\Omega/\mathrm{J}\Omega) = \mathrm{long}_{\mathrm{A}}(\Lambda/\mathrm{J})\) and the \(\kappa_{\mathrm{A}}\) -vector space \(\mathrm{Hom}_{\Lambda}(\kappa_{\Lambda}, \Omega/\mathrm{J}\Omega)\) has dimension 1.
\end{enumerate}

(i)⇒(ii): if Ω is dualizing, it is Macaulay and of dimension dim(A) (n° 1, prop. 1). The sequence x is A-regular since A is a Macaulay ring; by prop. 4, the A/J-module Ω/JΩ is dualizing, hence is a Matlis A/J-module (n° 1, example 1).

(ii)⇒(iii): under hypothesis (ii), one has \(\dim(\Omega)=\dim(A)\) and \(\dim(\Omega/J\Omega)=\dim(A/J)=0\) , so that the sequence x is secant for \(\Omega\) , hence \(\Omega\) -regular (§ 2, n° 3, th. 1).

(iii)⇒(i): under the hypotheses of (iii), the Λ/J-module Ω/JΩ is a Matlis A-module, hence is dualizing (n° 1, example 1); by prop. 4 the A-module Ω is dualizing.

\begin{enumerate}
\def\labelenumi{(\roman{enumi})}
\setcounter{enumi}{2}
\tightlist
\item
  \(\Leftrightarrow\) (iv): this follows from the remark of § 8, n° 3.
\end{enumerate}

